% ============================================================
% Digital Discovery Methods — honest structure-based benchmark (repositioned 2026-06-23)
% Numbers trace to data/final_model_metrics.json, structural_rebuild_metrics.json,
% dft_benchmark_comparison.csv, and the scripts in code/.
% ============================================================

\section{Methods}\label{sec:methods}

\subsection{Experimental target and dataset}\label{para:validation_hierarchy}

The prediction target is the \emph{experimental} singlet--triplet gap
$\Delta E_\text{ST}$. We extracted literature-reported values from a curated
corpus of TADF publications and matched them, by compound name, to molecules in
our descriptor set, retaining values in the physical range
\qtyrange{-0.3}{1.5}{\electronvolt}. Where a compound carried multiple independent
reports we used the median and recorded the inter-report spread; this scatter is
typically \qtyrange{0.1}{0.15}{\electronvolt} and bounds achievable accuracy
(\cref{subsec:ceilings}). The final benchmark comprises \num{231} donor--acceptor
molecules spanning \num{212} distinct Bemis--Murcko scaffolds~\cite{Bemis1996}.

The corpus does not record measurement medium (solvent, host) for individual
values, so the per-molecule median is a consensus over uncontrolled conditions
rather than a single, solvent-specific state. Because environmental electrostatics
and conformational disorder are recognised as primary obstacles to accurate TADF
modelling,\cite{Olivier2017} this consensus is best read as the central tendency of
experimentally accessible $S_1$--$T_1$ gaps; it trades physical specificity for
reduced variance. For the \num{102} molecules with repeat reports, a single random
report deviates from this median by \qty{0.072}{\electronvolt} on
average---comparable to the model error itself
(\cref{SI-tab:solvent})---which sets a floor on attainable accuracy independent of
model capacity.

We deliberately do \emph{not} train on the semi-empirical gap
$E_{S_1}-E_{T_1}$: because this quantity is, by construction, a difference of two
input features, regressing onto it measures arithmetic rather than learning
(\cref{subsec:circularity}).

\subsection{Cheap electronic-structure inputs and physics-informed descriptors}

Excited-state quantities used to build descriptors were obtained at the
semi-empirical level: GFN2-xTB~\cite{Bannwarth2019} geometries
(\texttt{xtb} 6.5.1) followed by sTDA~\cite{Grimme2017} for the lowest singlet
and triplet states. From the natural transition orbitals
(NTOs)~\cite{Martin2003} of each excitation we computed spatial overlap
descriptors that encode exchange-integral physics:
\begin{align}
  S_\text{he} &= \int \phi_\text{hole}(\mathbf{r})\,\phi_\text{electron}(\mathbf{r})\,d\mathbf{r}, \label{eq:overlap}\\
  \Delta\lambda_\text{A} &= |\lambda_\text{A}(S_1) - \lambda_\text{A}(T_1)|, \label{eq:delta_lambda}\\
  \text{Char}_\text{diff}^2 &= \Big[\textstyle\sum_i |q_i(S_1) - q_i(T_1)|\Big]^2, \label{eq:char_diff}\\
  \Delta S_\text{he} &= |S_\text{he}(S_1) - S_\text{he}(T_1)|. \label{eq:delta_overlap}
\end{align}
$S_\text{he}$ quantifies hole--electron overlap (related to the exchange
integral), $\Delta\lambda_\text{A}$ acceptor-participation differences,
$\text{Char}_\text{diff}^2$ charge-transfer character changes, and
$\Delta S_\text{he}$ the singlet--triplet overlap difference. The full set of
\num{35} NTO descriptors (the $S_1$/$T_1$ energy scalars excluded) is listed in
the Supporting Information; these descriptors require the sTDA excited-state
calculation, even though they omit the energies themselves.

\subsection{Feature sets and machine-learning models}

We compared three feature representations on the same molecules. Two require no
quantum chemistry, being computed from the 2D structure alone:
(i)~RDKit physicochemical descriptors~\cite{Landrum2016} and (ii)~Morgan
fingerprints (radius~2, \num{2048} bits). The third, (iii)~the NTO spatial
descriptors above, is derived from the sTDA excited state but excludes the energy
scalars. A reference set that additionally includes the sTDA $S_1$/$T_1$
energies was also evaluated. We used random-forest
regressors~\cite{Breiman2001} (\num{400} trees, scikit-learn~\cite{Pedregosa2011}),
which outperformed support-vector regression, ridge regression and a mean-value
baseline on this task.

\paragraph{Evaluation.}
All performance figures use Bemis--Murcko scaffold \texttt{GroupKFold}
($k=5$): every test molecule lies on a scaffold unseen in training, so reported
errors estimate generalisation to new chemotypes rather than interpolation.
We benchmark against a mean-value predictor and the raw sTDA gap, and report
\qty{95}{\percent} confidence intervals from \num{2000} bootstrap resamples of the
out-of-fold predictions.

\paragraph{Interpretability and reproducibility.}
Feature importance was quantified with exact TreeSHAP~\cite{Lundberg2017} on the
fitted forest. All scripts (\texttt{code/finalize\_model.py},
\texttt{code/rebuild\_structural\_model.py}) regenerate every reported number and
\cref{fig:model_performance}.

\paragraph{Bootstrap confidence intervals and power analysis.}
Point estimates of the rank correlation are complemented by
\num{5000}-replicate percentile bootstrap confidence intervals on the
out-of-fold predictions (\texttt{code/r2\_ci\_power\_analysis.py}).
For the headline model the Spearman $\rho = 0.36$ ($n = 231$) has a 95\%~CI of
\numrange{0.24}{0.47} ($p = \num{2.6e-8}$) and the Pearson $r = 0.50$ a 95\%~CI
of \numrange{0.26}{0.68}; both exclude zero, establishing statistically
significant association despite the wide $R^2$ interval.
A power analysis for the correlation test (Fisher $z$ transform, two-sided,
$\alpha = 0.05$) shows that detecting $\rho = 0.36$ at \qty{80}{\percent} power
requires only $n \approx 60$ (\qty{90}{\percent} power: $n \approx 79$); at
$n = 231$ the achieved power is $\approx 1.0$. The ranking-ability question is
thus amply powered, and the wide $R^2$ interval is \emph{not} a sample-size
problem: it reflects the high variance of a variance-explained statistic when a
handful of high-gap outliers dominate the residual sum of squares. We therefore
report Spearman~$\rho$ as the primary metric for triage utility, with the full
set of metrics and intervals collected in \cref{SI-tab:metrics_ci}.

\subsection{Reference-level benchmark}

To bound the meaning of any model accuracy, we compared the computed reference
across functionals for \num{7} representative emitters in gas and toluene,
contrasting B3LYP/6-31G(d) with CAM-B3LYP/def2-TZVP (Gaussian~16~\cite{Gaussian16};
ORCA~6.1.0~\cite{Neese2020}). The two levels differ in $\Delta E_\text{ST}$ by a
mean absolute \qty{0.30}{\electronvolt} (up to \qty{0.58}{\electronvolt}, with
cases of reversed relative magnitude); full values are tabulated in
\cref{SI-tab:dft_benchmark}. This functional dependence, together with
experimental inter-report scatter, sets the accuracy ceiling discussed in
\cref{subsec:ceilings}. The same CAM-B3LYP/def2-TZVP calculations additionally
yielded excited-state-optimised (adiabatic) gaps and the two lowest triplets
($T_1$, $T_2$) for \num{14} emitters, used to bound the vertical approximation
(\cref{subsec:ceilings}) and the triplet-manifold scope (\cref{subsec:limitations});
see SI.

\subsection{Heuristic enumeration of an extended library}\label{sec:library_expansion}

To explore chemical space beyond the annotated set, we enumerated \num{22194}
donor--acceptor combinations from literature-validated fragments (carbazole,
phenoxazine, acridine donors; triazine, benzonitrile, sulfone acceptors) under
simple constraints (molecular weight $<\qty{900}{\dalton}$, $\geq 2$ aromatic
rings, $\geq 1$ N/O/S heteroatom). These were ranked by a rule-based TADF score
\begin{equation}
  \text{TADF score} = 0.35\,f_\text{MW} + 0.25\,f_\text{LogP} + 0.20\,f_\text{rings} + 0.20\,f_\text{hetero},
  \label{eq:tadf_score}
\end{equation}
with each $f_i\in[0,1]$ a normalised feature score based on established design
heuristics~\cite{Liu2012,Uoyama2012,Wong2017}. The score is a coarse
synthesis-planning filter, \emph{not} a calibrated $\Delta E_\text{ST}$ predictor;
the enumerated molecules are not quantum-validated and are reported as leads only.

\subsection{Data and code availability}

The experimental corpus, descriptor tables, trained models, and all analysis and
figure-generation scripts are openly available on Zenodo
(DOI: 10.5281/zenodo.17436069) under a CC-BY-4.0 licence, so that every number and
figure in this work can be reproduced and the benchmark extended.
